% Part: first-order-logic
% Chapter: syntax-and-semantics
% Section: subformulas

\documentclass[../../../include/open-logic-section]{subfiles}

\begin{document}

\olfileid{fol}{syn}{sbf}

\olsection{\printtoken{P}{subformula}}

\begin{explain}
نحتاج كثيرًا إلى النظر في ال!!{formula}s التي تألّفت منها
!!a{formula} معطاة؛ وتلك هي \emph{!!{subformula}s} لها. وكل
!!{formula} تُعدّ !!a{subformula} من نفسها. أما الصيغة الجزئية من~$!A$
إذا كانت غير~$!A$ نفسها، فتسمى \emph{!!{subformula} حقيقية}.
\end{explain}

\begin{defn}[!!^{subformula} مباشرة]
إذا كانت $!A$ !!a{formula}، فتُعرّف \emph{ال!!{subformula}s
المباشرة} لـ~$!A$ استقرائيًا كما يأتي:
\begin{enumerate}
\item ليس لل!!{formula}s الذرية !!{subformula}s مباشرة.

\tagitem{prvNot}{\indcase{!A}{\lnot !B}{ال!!{subformula} المباشرة
    الوحيدة لـ~$\indfrm$ هي~$!B$.}}{}

\item \indcase{!A}{(!B \ast !C)}{ال!!{subformula}s المباشرة
  لـ~$\indfrm$ هي $!B$ و$!C$ (حيث $\ast$ أي رابط ثنائي).}

\tagitem{prvAll}{\indcase{!A}{\lforall[x][!B]}{ال!!{subformula}
    المباشرة الوحيدة لـ~$\indfrm$ هي~$!B$.}}{}

\tagitem{prvEx}{\indcase{!A}{\lexists[x][!B]}{ال!!{subformula}
    المباشرة الوحيدة لـ~$\indfrm$ هي~$!B$.}}{}
\end{enumerate}
\end{defn}

\begin{defn}[!!^{subformula} حقيقية]
إذا كانت $!A$ !!a{formula}، فتُعرّف \emph{ال!!{subformula}s
الحقيقية} لـ~$!A$ عوديًا كما يأتي:
\begin{enumerate}
\item ليس لل!!{formula}s الذرية !!{subformula}s حقيقية.

\tagitem{prvNot}{\indcase{!A}{\lnot !B}{ال!!{subformula}s الحقيقية
    لـ~$\indfrm$ هي~$!B$ مع جميع ال!!{subformula}s الحقيقية
    لـ~$!B$.}}{}

\item \indcase{!A}{(!B \ast !C)}{ال!!{subformula}s الحقيقية
  لـ~$\indfrm$ هي $!B$ و$!C$، مع جميع ال!!{subformula}s الحقيقية
  لـ~$!B$ وتلك الخاصة بـ~$!C$.}

\tagitem{prvAll}{\indcase{!A}{\lforall[x][!B]}{ال!!{subformula}s
    الحقيقية لـ~$\indfrm$ هي~$!B$ مع جميع ال!!{subformula}s الحقيقية
    لـ~$!B$.}}{}

\tagitem{prvEx}{\indcase{!A}{\lexists[x][!B]}{ال!!{subformula}s
    الحقيقية لـ~$\indfrm$ هي~$!B$ مع جميع ال!!{subformula}s الحقيقية
    لـ~$!B$.}}{}
\end{enumerate}
\end{defn}

\begin{defn}[!!^{subformula}]
ال!!{subformula}s لـ~$!A$ هي $!A$ نفسها مع جميع ال!!{subformula}s
الحقيقية لها.
\end{defn}

\begin{explain}
وبين تعريف ال!!{subformula}s المباشرة وتعريف ال!!{subformula}s
الحقيقية فرق ينبغي تبيّنه. فقد عيّنّا في الأول ال!!{subformula}s
المباشرة للصيغة~$!A$ بحسب كل صورة يجوز أن تكون عليها~$!A$، فكان
تعريفًا صريحًا بالحالات، على ترتيب التعريف الاستقرائي لمجموعة
ال!!{formula}s. وسلكنا في الثاني ترتيب تعريف مجموعة
ال!!{formula}s أيضًا، غير أننا أدخلنا في كل حالة ال!!{subformula}s
الحقيقية لل!!{formula}s الأصغر $!B$ و$!C$، إلى جانب هاتين
ال!!{formula}s نفسيهما. وبذلك صار التعريف \emph{عوديًا}.
وضابط العود في تعريف دالة على مجموعة محدّدة استقرائيًا، وهي هنا
ال!!{formula}s، أن تستعمل بعض حالات التعريف الدالة نفسها.
ولا يستقيم ذلك إلا إذا اقتصر الاستعمال على قيم الدالة عند وسائط
سابقة للوسيط الجاري تعريف قيمته. فحين نحدّد «!!{subformula} حقيقية»
لـ~$(!B \ast !C)$، إنما نأخذ ال!!{subformula}s الحقيقية
لل!!{formula}s السابقة في البناء، أعني $!B$ و$!C$.
\end{explain}

\begin{prop}
\ollabel{prop:subfrm-trans}
إذا كانت $!B$ صيغة جزئية من~$!A$، وكانت $!C$ صيغة جزئية من~$!B$،
لزم أن تكون $!C$ صيغة جزئية من~$!A$؛ أي إن علاقة الصيغة
الجزئية علاقة متعدية.
\end{prop}

\begin{prob}
أثبت \olref[fol][syn][sbf]{prop:subfrm-trans}.
\end{prob}

\begin{prop}
\ollabel{prop:count-subfrms}
لتكن $!A$ صيغة مجموع ما فيها من الروابط والمكممات $\OLId{latin-ordinary}{n}$.
فعدد الصيغ الجزئية لـ~$!A$ لا يتجاوز $\OLMathNumeral{2}\OLId{latin-ordinary}{n}+\OLMathNumeral{1}$.
\end{prop}

\begin{prob}
أثبت \olref[fol][syn][sbf]{prop:count-subfrms}.
\end{prob}

\end{document}
